% Discussion paragraph carrying the core of Storyline 2 (agreement from private revision), for either storyline.
% In the manuscript, refer to Fig.~\ref{fig:recognition} for the Shunning example.
Our results also bear on how agreement about reputations arises. Agreement is usually attributed to mechanisms that align information across observers, such as institutions, gossip and opinion synchronization\cite{radzvilavicius2021,kawakatsu2024,murase2024synchronization}, or to observers who draw on several observations at once\cite{michelmata2024}. In our model, observers share nothing and update on one observation at a time, yet they largely agree. Two independent observers who each regard an individual as good with probability $r$ disagree with probability $2r(1-r)$, so agreement follows whenever private judgements become reliable. Under the punishment dilemma, reliability is not automatic, because the assessments an individual receives depend on what observers think of its recipients. Inertia produces it when it is placed in the direction that separates the assessments of cooperators from those of defectors. Otherwise observers can agree on the wrong verdict, as under one-step Shunning, where they regard nearly everyone as bad (Fig.~2). Agreement through private revision complements information sharing rather than replacing it. It remains available where shared information is costly or untrusted, as in polarized societies\cite{iyengar2019,michelmata2024}, and it can be combined with institutions or gossip where they exist.
